\begin{abstract}
Knowledge tracing (KT) underpins adaptive learning engines, yet most KT
research targets the mathematics domain, leaving the
programming/algorithm-learning domain comparatively under-studied.
Furthermore, KT work overwhelmingly optimises next-step prediction
accuracy (AUC), while the link between predictive accuracy and the
quality of the downstream instructional decision---which exercise to
give next---remains weakly established, and is largely unexamined for
programming. This work investigates, for programming-concept learning,
whether KT models that exploit programming-specific signals
(knowledge-component tags and code-structure features derived from
abstract syntax trees) (a)~predict learner mastery more accurately than
response-only models, and (b)~whether any accuracy gains translate into
better next-exercise sequencing. We build a reproducible programming-KT
evaluation pipeline over multiple public datasets (CodeWorkout,
FalconCode) with standard baselines (DKT, SAKT, AKT, simpleKT),
implement a knowledge-component / code-structure-aware variant on top of
a strong baseline, and evaluate both predictive accuracy (AUC, RMSE) and
downstream sequencing utility via offline policy evaluation. We further
analyse the correlation between predictive accuracy and sequencing
utility, and assess cross-dataset and cross-domain (programming vs.
mathematics) generalisation.
% \todo{Fill in headline quantitative results once experiments are done.}
\end{abstract}
